\scanpage{35}
\uline{Définition} 2.5. Soient \(X,Y\ldots\) et considérons une application
\diagram{\begin{tikzcd}[ampersand replacement=\&]
\varphi:\uline{Y(k)}\arrow[r]\&|[draw,ellipse]|A^i_{\mathrm{alb}}(X)
\end{tikzcd}}
On dit que cette application est \uline{algébrique} si

1°) L'hom \(\deleted{\unsure{Z}}\operatorname{Cyc}_0(Y)\longrightarrow A^i_{\mathrm{alb}}(X)\) correspondant se factorise pr \(\deleted{A}A_{0\,\mathrm{alb}}(Y)\simeq A_Y^*(k)\), i.e. \deleted{\unsure{elle}} est nulle sur les \(\nu\) de degré 0 tels que \figure{\begin{tikzpicture}[baseline=(sy.base)]\node[draw,ellipse,inner sep=2pt] (sy) {$S_Y(\nu)=0$};\end{tikzpicture}}.
\margin{\diagram{\begin{tikzcd}[ampersand replacement=\&]Y\arrow[r]\&A_Y^1\arrow[r,"N"]\&A_Y^N\end{tikzcd}}}

2°) L'hom induit
\diagram{\begin{tikzcd}[ampersand replacement=\&]
|[draw,ellipse,label={left:$\varphi:\,$}]|A_Y^{\unsure{0|*}}(k)\arrow[r]\&A^i_{\mathrm{alb}}(X)
\end{tikzcd}}
est tel que pr \(N>0\) convenable, \(\underline{N\varphi}(x)=\varphi(Nx)\) pourra être induit pr un cycle algébrique
\[
Z\in A^i_{\mathrm{alb}}(X\times A_Y)
\]

\uline{Prop.} 2.6. Si \(\varphi\) est définie pr un \(Z\in A^i_{\mathrm{alb}}(X\times Y)\), elle est \uline{algébrique}.

C'est 2.1.
\margin{Va être généralisé en 2.7. (ii)}

\uline{Prop.} 2.7. \inserted{(0) Forment un \uline{\unsure{k-groupe|s-groupe}}. Si \(N\varphi\) est alg., \(\varphi\) l'est.}

(i) Fonctorialité en \(Y\). \deleted{Alors \unsure{\(\varphi\) induit}}

(ii) Fonctorialité en \(X\), plus gén., si on a un hom \(A^i_{\mathrm{alb}}(X)\longrightarrow A^{\unsure{j}}_{\mathrm{alb}}(X')\) défini pr une classe de correspondances algébriques \uline{à coeff.} \deleted{\uline{\unsure{rationnels}}} entiers.
\mspage{34}
